%!TEX root = ../dissertation.tex

\chapter{Models}
\label{chp:models}

This chapter discusses the experimental setup for the predictive models used for testing the explainability framework. It begins by outlining the datasets used as well as the processing that was applied to them. Afterwards, an overview of the pre-processing done is provided to ensure the chosen predictive model was working with the best possible hyperparameters and the comparison made against other baseline models to validate the chosen architecture. Finally, a rundown of the explainability framework that was applied to the predictive model is provided.

\section{Use Cases}
    \subsection{Order Management}
        \label{sec:dataset_order_management}
    		The first of the two datasets used for validating the proposed explainability framework was a synthetic OCEL 2.0 benchmark dataset named Order Management \cite{knopp2023a}. The general process pipeline in the dataset is as follows; activities are placed in parentheses at relevant points of the explanation.
    
    		The process begins when a customer places an order for different products (Place Order). Each product has a price (which can vary on different dates) and a weight. When a customer places an order it is assigned to an employee to act as a sales-representative. The employees are tasked with confirming the customer's order (Confirm Order) as well as processing the payment (Pay Order) and reminding customers of payments that have not been received yet (Payment Reminder). 
    		Running parallel to this is the company's shipping process. After an order is received the order's items (unique identifiers for the products in each order) are checked against the company's stock. If an item is unavailable it is marked as out of stock (Item Out of Stock) and reordered by the warehouse employee (Reorder Item). Items that are ready for delivery are collected (Pick Item) for placement into packages that are addressed to a single customer. A package may contain items related to multiple orders, and an order's volume may be distributed over multiple packages.
    		After the readied items are collected into a package (Create Package) the package is picked up by a shipment employee for transport (Send Package).
    		Finally, the package is shipped. Deliveries may fail repeatedly (Failed Delivery) until a successful delivery is achieved (Package Delivered).
    
    		The KPIs chosen to analyze in this dataset were Elapsed time from Orders to PackageDelivered, which identifies the time from the creation of the order until the delivery of the final package related to that order, and Elapsed time from Orders to PayOrder, which identifies the time from the creation of the order until the order is paid by the customer.
    
    		This dataset is composed of few, well-connected items so a model using only 2 layers was sufficient to create a path from the viewpoint to any other node type. This echoes the design decision of Smit et al., who used a similar model design, although their work was based around OCEL version 1.0 of this dataset \cite{smit2024}.
            
    \begin{table}[t!]
        \centering
        \caption{Summary of the datasets used in this thesis.}
        \label{tab:dataset_summary}
        \begin{tabular}{lrrrrrr}
            \toprule
            \multirow{2}{*}{OCEL} & \multicolumn{2}{c}{Events} & \multicolumn{3}{c}{Objects} & \multirow{2}{*}{Traces}\\
            \cmidrule(lr){2-3} \cmidrule(lr){4-6}
             & Instances & Types & Instances & Types & Attributes &  \\
            \midrule
            Order Management & 20801 & 11 & 10815 & 6 & 7 & 2000\\
            Logistics & 121913 & 14 & 47591 & 7 & 7 & 2000\\
            \bottomrule
        \end{tabular}
        \end{table}
    
	\subsection{Logistics}
    \label{sec:dataset_logistics}
		The second dataset used to validate the proposed explainability framework is a different synthetic OCEL 2.0 benchmark dataset named Logistics \cite{knopp2023b}. The general process pipeline in the dataset is as follows; activities are placed in parentheses at relevant points of the explanation.

		The process begins when a customer's order is registered at the company (Register Customer Order). After registration, a transport document in which the details of the rest of the process will be recorded is created (Create Transport Document). 
		A logistics service provider is contacted to coordinate the transport of the ordered goods. Twice a week, the service provider sends a vehicle to a terminal to transport containers of ordered goods from the terminal to the seaport. Since the company depicted is not the only company that uses the provider's services available capacities may vary from vehicle to vehicle. Once the logistics provider receives the transport document, it books vehicles according to availability and priority (Book Vehicles). Once the dates for the transport of the goods are set, a container depot is contacted to reserve the required containers (Order Empty Containers).
		When a container vehicle's departure approaches the required goods are collected (Collect Goods). A truck is sent to the container depot to pick up the empty containers (Pick Up Empty Container). The truck is loaded with the collected goods (Load Truck), and it drives the full container to the terminal (Drive to Terminal).
		At the terminal, the container is picked up by a forklift and weighed (Weigh). The container may then be placed in the storage location at the terminal to await departure (Place in Stock). Finally, the container is taken to the loading bay and then moved onto the vehicle (Bring to Loading Bay, Load to Vehicle) where it departs at a set time (Depart).
		A container may miss a vehicle's departure and be rescheduled for the next possible vehicle (Reschedule Container).

		The KPIs chosen to analyze in this dataset were Elapsed time from CustomerOrder to Depart, which identifies the time from the creation of the customer's order until the departure of the final container associated with it, and Elapsed time from CustomerOrder to LoadToVehicle, which identifies the time from the creation of the customer's order until the container is loaded onto the vehicle.

		This dataset differs from Order Management in its structure, with a higher quantity of object types and connections represented. Given this added complexity it was necessary to increase the chosen model's depth beyond 2 layers. This is due to the fact that a model at only 2 layers was incapable of creating a connection between the viewpoint node chosen and all the other node types present, causing the model to collapse to a near-constant predictor regardless of training budget. In order to find the proper layer depth the model was evaluated at 3, 4, and 5 layers for each KPI; layer depth 3 showed the best performance for Elapsed time from CustomerOrder to Depart, while Elapsed time from CustomerOrder to LoadToVehicle instead settled on layer depth 4.

	\subsection{Event log processing}
    \label{sec:dataset_data_processing}

	This section expands upon the OCEL concepts introduced in Section 2.1.1 and formalized in Section \ref{sec:dataset_OCEL}. As stated there, the OCEL is generated by identifying the elements present in a process and the relationships that appear between them as events in the process are executed. In the case of the datasets the information related to objects, events and their relationships is found in an SQLite file containing a relational database of tables. These tables can be split into three main groups: tables relating to events, tables relating to objects and tables detailing the relationship between these elements. Table \ref{tab:dataset_summary} provides a summary of the number of each instance present in each dataset. Tables relating to events provide a row by row accounting of all events of a certain type that occurred and the timestamp of the event's occurrence. Tables relating to objects provide similar information. For each object type they provide an OCEL identifier, a timestamp representing the time the object was created, a column indicating a change in one of the object table's columns, as well as columns for each attribute associated with the given object. For example, table object\textunderscore packages contains information on each package type object in the database. This table consists of an identification column, an ocel\textunderscore time column, an ocel\textunderscore changed\textunderscore field column, and a weight column, that provides the weight of the package.
    
    Following traditional relational database conventions, the relationships between the elements in these tables are kept in separate tables labeled event\textunderscore object and object\textunderscore object (for event-to-object and object-to-object relationships respectively). 
    All these elements can be extracted from the original database through SQL queries to generate the OCEL. Following the formalizations provided in Section \ref{sec:dataset_OCEL} the OCEL creation process is as follows.
    A viewpoint object is selected from among the object types present in the database (for example, Order). For each object of the chosen type a sequence of related objects are identified through the use of the object\textunderscore object table. Once all related objects are identified each event related to the objects is found by querying the event\textunderscore object table. The resulting collection of events is then ordered according to their associated timestamp to create a basic event log. The log is then expanded upon by associating each related object of each event with its attributes, thus creating the OCEL.

    It is important to note that while event-to-object and object-to-object relationships are directly provided by the database event-to-event relationships need to be inferred from the data. This process consists of identifying events with related objects that have a 1-hop distance between each other.

    While the database provides a set of attributes for each object node, event nodes tend to have no related attributes. In order to enhance the information provided to the predictive model each node is feature-enhanced following the work of Adams et al. \cite{adams2022} by calculating object-centric features using information that is present in the data but not explicitly stated. Such features may include values like waiting time which measures how long an event node had to wait before appearing in the process execution's graph, or an aggregate count of object types that have appeared so far in the process execution. 

    Given the wide array of values that can appear in a given node's features it is important to perform a process of standardization to ensure that large scale values, such as price, do not get assigned greater importance by the model when compared to smaller scale values such as weight. This standardization is also performed on the target value to facilitate faster convergence on the GNN's gradient descent-based algorithm. Standardization is performed using Z-score normalization. This process consists of calculating the mean and standard deviation of each feature value in the training set, then subtracting the mean from the value of each feature and then dividing by the standard deviation. This process ensures that each value is transformed to have a mean of 0 and a standard deviation of 1.

    After these processes have been completed, the graph representation of each process execution and its associated prefixes is created from the OCEL's data. This results in the dataset being a series of subgraphs, as opposed to a single global schema of interconnected elements. Future work may follow a global approach, similar to the one proposed by de Leoni \& Volpato \cite{deleonieVolpato}, to see if it provides better results, but that implementation is outside the scope of this thesis.


\section{Heterogeneous Graph Neural Network for Predictive Process Monitoring}
    In order to apply the proposed explainability layer it was first necessary to train a predictive model that could accurately predict values using the heterogeneous dataset that is created following the procedures outlined in Section \ref{sec:dataset_OCEL}. This process began by defining a procedure for performing a temporal train/test split of the data. Afterwards, a model architecture is chosen that is well suited for the heterogeneous graph structure the dataset is encoded in. By using a hyperparameter sweep it is possible to identify the best hyperparameters to use in the chosen model architecture and then perform a training and validation loop. Finally, the results are checked against baseline models to verify whether the chosen implementation is comparably effective against versions using other encodings, such as a homogeneous GNN or a traditional gradient boosted model on tabular data.

    \subsection{Train/test split of the event logs}
        The set of process executions and associated prefixes of the datasets all have an important temporal aspect related to them. The events in the processes follow a strict chronological order, indicating that some processes are completed long before other processes even begin. Therefore, it is important to ensure that there is no information leakage from future events into past ones when training the model. To ensure this, the train/test split of the data is performed ensuring the temporal separation between training, validation and test sets following the procedures used by de Leoni \& Volpato\cite{deleonieVolpato}.
        In order to achieve this the universe of process executions was ordered chronologically according to the timestamp of their final associated event. The first 70\% of the total process executions were selected for use in the training set. Any process executions outside this 70\% that were considered active at the time of the completion of the final training set event (i.e. the process had begun before the completion of the final training set event and will finalize after the completion of the final set event) were eliminated from the dataset to ensure that all the process executions present in the other sets were temporally separated from the training data. The process is repeated on the remaining process executions: the first 50\% of the remaining ordered process executions are selected for the validation set, and all the process executions that remain after ensuring temporal separation again are assigned to the test set. This method inevitably results in process executions that are not taken into account, with the real train/test split percentages not equaling 100\%. However, this is offset by the assurance of no information leakage between sets.

    \subsection{Model architecture choice and evaluation metrics}
        Two possible heterogeneous model architectures were considered for this work. The first is Smit et al.'s \cite{smit2024} k-dimensional GNN. The other is Zhai et al.'s \cite{zhai2025} HGT implementation. Since this work is adapting Zhai et al.'s explainability layer, it was decided to also work with their proposed heterogeneous model architecture; this way it could also offer a point of comparison against Smit et al.'s implementation.
        A major difference in both architectures is the way they handle the heterogeneous aspect of the input data. The k-dimensional GNN model transforms a homogeneous model design into a heterogeneous version through the use of PyTorch Geometric's to\textunderscore hetero transformer. This transformation is achieved by duplicating the homogeneous model's message passing instance for each node type that exists in the heterogeneous input provided to the model. The chosen architecture, on the other hand, follows a strictly heterogeneous architecture built around heterogeneous graph transformers which are designed with per-type projection matrices and cross-type attention from the ground up. 

        To evaluate the performance of the chosen regression model two metrics were employed: Mean Absolute Error (MAE) and Root Mean Square Error (RMSE). These are widely recognized metrics for assessing the accuracy and predictive quality of regression models. Here, a definition of each metric is provided, where $n$ is the number of observations in the test set, $y_i$ represents actual values and $\hat{y}_i$ represents predicted values.

        \begin{itemize}
            \item Mean Absolute Error (MAE): MAE quantifies the average absolute difference between the predicted values and the actual values. This measure provides insight into the overall error magnitude such that a lower MAE indicates smaller errors, and therefore better model performance.\\
            \begin{center}
            $MAE = \frac{1}{n}\displaystyle\sum_{i=1}^{n}|y_i - \hat{y}_i|$
            \end{center}
            
            \item Root Mean Squared Error (RMSE): RMSE is a quadratic scoring rule that calculates the square root of the average of squared differences between the predicted and actual values. Same as MAE, lower RMSE values indicate better model performance and provide complementary insight by highlighting larger errors.\\
            \begin{center}
            $RMSE=\sqrt{\frac{1}{n} \displaystyle\sum_{i=1}^{n}(y_i-\hat{y}_i)^2}$
            \end{center}
    
            \item Coefficient of Determination (R-Squared): $R^2$ is a statistical measure that indicates the proportion of the variance in the dependent variable that is predictable from the independent variables. It complements the other metrics by also providing a measure of how well the model handles the variability in the data. A higher value of $R^2$ indicates better model performance.\\
            \begin{center}
            $R^2 = 1-\frac{\sum_{i=1}^{n}(y_i-\hat{y}_i)^2}{\sum_{i=1}^{n}(y_i-\overline{y})^2}$ 
            where $\overline{y}$ is the mean of the actual values $y_i$
            \end{center}
        \end{itemize}
        
    \subsection{Hyperparameter sweep}
        \begin{table}[t!]
            \centering
            \caption{Final chosen hyperparameters for this thesis's two datasets. Logistics is split into its two KPIs since, unlike Order Management, they do not share the same final configuration.}
            \label{tab:final_hyperparameters}
            \begin{tabular}{lrrr}
            \toprule
            Hyperparameter & \shortstack{Order\\Management} & \shortstack{Logistics\\(Depart)} & \shortstack{Logistics\\(LoadToVehicle)} \\
            \midrule
            Number of epochs & 30 & 100 & 100 \\
            Patience & 4 & 10 & 10 \\
            Number of message-passing layers & 2 & 3 & 4 \\
            Activation function & PReLU & PReLU & PReLU \\
            Number of heads & 2 & 2 & 2 \\
            Optimizer & Adam & AdamW & AdamW \\
            Number of hidden dimensions & 64 & 16 & 48 \\
            Learning rate & 0.001 & 0.005 & 0.003 \\
            Weight decay & 0.0 & 1e-05 & 1e-05 \\
            \bottomrule
            \end{tabular}
            \vspace{2pt}
        \end{table}

        Having chosen the model architecture and the proper metrics for assessing the model's performance it was necessary to identify the most effective hyperparameter configuration per dataset and chosen KPI. In order to do this a grid search was performed over the model's number of hidden channels, and its learning rate. During this search it was found that a mismatch between the weight decay used during the sweep (0.001) and the weight decay used in the actual regression testing (0.00001) led to the zeroing out of several node types' input projections entirely. The bug was identified and fixed during testing.

        Aside from the values already mentioned the logistics dataset, given the problem discussed in Section \ref{sec:dataset_logistics}, also searched for an optimal number of layers greater than 2, specifically among the values 3, 4, and 5. It's worth mentioning that any model that used only 1 layer for its architecture made any node type without a direct edge to the viewpoint unreachable for the model. The values used as possible choices for the search are informed by priors taken from Smit et al.'s \cite{smit2024} and de Leoni \& Volpato's \cite{deleonieVolpato} reported findings on which values worked best for each of the chosen datasets. 


    \subsection{Training and validation}
        With the model's hyperparameters set, the training and testing the model's predictions was able to begin. An ADAM optimizer (AdamW for the two logistics KPIs, as detailed in Table \ref{tab:final_hyperparameters}) is used, combined with an L1Loss function. In order to provide an accurate point of comparison to Smit et al.'s \cite{smit2024} work the testing parameters are aligned to the ones they reported on by setting up a testing environment with 30 epochs and a patience value of 4 epochs (i.e. if the model has not improved its prediction in 4 epochs the training loop ends). 
        On the other hand, the logistics dataset is not considered in Smit et al.'s work so the epoch and patience values did not need to align with any past results. Additionally, it was found that at the selected patience level of 4 the predictive model undertrains. To combat this, the training epochs were increased from 30 to 100 and its patience value from 4 to 10.

    \subsection{Baselines for comparison}
    \label{sec:models_baselines}
        In order to verify if the chosen heterogeneous architecture is truly more effective at performing predictions than past encodings a comparison was set up between the model and three baselines: a mean predictor used as a lower bound, Gradient Boosted Trees on a tabular encoding and a homogeneous GNN trained on a homogeneous representation of the dataset. These proposed baselines are adapted from Smit et al.'s work \cite{smit2024}.

        The homogeneous GNN model uses a Graph Convolutional Network (GCN) \cite{kipf2017} architecture with ReLU activation. Rather than being independently hyperparameter-tuned, its number of layers and hidden dimension are reused from the heterogeneous model's own tuned values for the given dataset and KPI (Table \ref{tab:final_hyperparameters}), keeping the two architectures at a comparable scale. After message passing, the resulting node embeddings are aggregated into a single graph-level representation through mean pooling, which is then passed through a two-layer fully-connected regression head to produce the final scalar KPI prediction. The homogeneous representation of the input data is constructed as a graph with only the event nodes of the original heterogeneous graph and their relationships present. Each event node has all the features present in its heterogeneous representation as well as a target value taken from the viewpoint object node, thus allowing for a direct comparison between a feature-enhanced heterogeneous encoding versus a feature-enhanced homogeneous encoding.

        Additionally, the heterogeneous encoding was checked to see if it provides a benefit for predictions over longer prefixes. The hypothesis was that the heterogeneous models would outperform the baselines significantly at greater prefix depths, becoming more precise the further along the process is. This is especially important for real world applications, as the predictions given by a model on a recently started process should be less accurate than the prediction it can provide on a much more mature one with more information to draw from.

\section{Explainability Framework Design}
    As outlined in Section \ref{sec:rw_xai} the proposed explainability framework is composed of four distinct components meant to enhance an end user's understanding of the predictive model, the reasons behind its predictions, and the ways in which those predictions can be changed through changes in the input data.

    \subsection{Perturbation-based explanation (LOO + Shapley Values)}
        The first component of the explainability framework is a perturbation-based analysis performed using a combination of Leave-One-Out (LOO) methodology and Shapley Values. As explained in Section \ref{sec:dataset_perturb} LOO is used to identify the most impactful features in the input graph in order to reduce the explanation space where the full Shapley Values method needs to be applied. The stakeholder is given the ability to decide on the $K$ number of most impactful elements they wish to view from the graph. Those elements are identified and then their impact is fully quantified with the Shapley Values method. 
        The stakeholder is then presented with a series of bar graphs identifying the most important elements in a given process execution and their calculated impact on the prediction should they be removed as well as a table summary of the same information to ensure its proper understanding. Additionally, an explanation subgraph is constructed using the identified elements and also presented to the stakeholder.
        The ranking of node-importance obtained through this explanation method is also shown in comparison to GNNExplainer's ranking as a validation check. As outlined in Section \ref{sec:dataset_perturb} GNNExplainer's learned mask is a [0,1]-bounded attention weight rather than a physically meaningful prediction-magnitude shift, and its heterogeneous implementation can't provide edge importance on a natively heterogeneous model like the chosen HGT architecture. Therefore, its role is limited to serving as only a comparison point.

        This analysis can also be performed at an aggregate level, quantifying the average impact of each feature found in a collection of process executions as opposed to a single trace. The aggregate values are obtained by calculating the feature's impact on its respective trace then averaging it over the number of appearances of that particular feature. The stakeholder is shown the resulting analysis in the form of a bar graph and a table similar to the one used to present the single trace results in order to maintain consistency in the presentation.

        It is important to note that the zero-vector substitution the LOO method uses to mask an element was itself validated against two other options: substituting the removed value with NaN or using a continuous soft mask like GNNExplainer. The first option was found to be untenable since using NaN led to a break in HGTConv's attention-weighted propagation. On the other hand, the second approach confirmed on a representative case that the current hard-zero design already is the fully-suppressed limit of a GNNExplainer-style soft mask.

    \subsection{Gradient/feature-based explanation (Gradient ⋅ Input and Integrated Gradients)}
        The second component of the explainability framework is used to complement the perturbation-based analysis of the first component by providing a finer-grained view of which features drive the important elements that component highlights. As explained in Section \ref{sec:dataset_gradient} this is done through the use of gradient-based methods applied across the model.
        At trace-level this analysis provides the stakeholder with a bar graph showing the attribution values for the top $K$ attributes in the current trace that are impacting the prediction. While the attribution value gives a general understanding of the importance of each feature it carries no dimensional value so it may be difficult to interpret. In order to improve the stakeholder's understanding each feature selected through a gradient method is then used to perform a Shapley Value analysis in order to provide the stakeholder with an interpretable magnitude of the shift the feature can have on the prediction. This attribution can be performed with either Gradient ⋅ Input or Integrated Gradients to allow for validation of each method. 

        Similar to how the perturbation-based explanations are shown, at an aggregate level the component analyzes the average impact of each feature of every node of a given type in the test set over their respective model predictions. The identified elements are shown to the user, ranked by importance and side by side with a bar graph showing the quantified impact on the prediction obtained through aggregate Shapley Value analysis on those features.


    \subsection{Counterfactual explanation}
    \label{sec:model_counterfactual}
        The third component of the framework is the counterfactual explanation which is used to provide the user with a direct comparison between existing subgraphs in the dataset that are similar but have different predicted values. As formalized in Section \ref{sec:dataset_cf} the analysis is broken into two parts: first, selecting the set of counterfactual candidates to compare against the factual graph and second, calculating the dissimilarity score between the factual graph and each possible candidate to identify the most similar graph with the desired prediction value.

        The first part of the process adapts the work of Huang et al. \cite{huang2022} and their LORELEY model for predictive processes. Their work outlines a way to identify possible candidate subgraphs within the universe of process executions, however, their work is designed for a classification problem. Three steps were taken to adapt their work for this thesis's regression setting. First, a minimum-meaningful-difference threshold is set instead of a simple class filter. This way the end user can select the minimum value difference they would like to see between the subgraphs when comparing them. Second, a set of explicit prefix-depth-matching rules were applied so candidates are compared at similar points in their process lifetime rather than across arbitrary lengths. Finally, the end user is allowed to decide whether the candidate values should be lower or greater than the current graph's prediction, this way the process can be extended for use cases such as revenue analysis where a higher predicted value is the goal.

        The second part adapts the dissimilarity metric used by Zhai et al \cite{zhai2025} to calculate how similar each of the chosen candidate graphs are to the factual graph. After calculating the dissimilarity metric for each pair the candidate with the lowest dissimilarity metric (i.e. the most similar graph) is chosen as the counterfactual example to show to the stakeholder.

        The stakeholder is then shown a graph representation of the factual graph as well as the chosen counterfactual. They're also shown a bar chart with the values of the four dissimilarity components that make up the final dissimilarity metric so they're able to identify which elements are driving the difference. Finally, the end user is provided with the specific elements that differ between each graph. With all these elements the stakeholder is able to identify the element driving the difference.

    \subsection{Fidelity metrics}
        The stakeholder is also presented with the fidelity metrics defined in Section \ref{sec:dataset_metrics}, calculated for the perturbation-based explanations and the resulting subgraph. By reviewing the fidelity values, the resulting characterization value as well as the sparsity of the explanation the stakeholder is able to judge how effective the proposed explanation is at assessing the model's functionality \cite{stevens2022}. 
